\section{정간보 텍스트 문법}
\label{sec:grammar}

\subsection{설계 목표}

문법 설계의 목표는 네 가지다.
(1) \textbf{가독성}: 국악 전공자가 문서만 보고 5분 안에 익힐 수 있을 것.
(2) \textbf{구조 동형성}: 텍스트의 시각적 배열이 악보 지면의 격자 구조와
일대일로 대응할 것 --- 줄이 곧 각이고, 구분자가 곧 정간선이다.
(3) \textbf{완전성}: 율명·옥타브·분박·붙임·이음·쉼표·숨표·시김새와
그 미세 조정값까지 지면에 나타나는 모든 정보를 잃지 않고 담을 것.
(4) \textbf{도구 친화성}: 일반 텍스트이므로 복사·붙여넣기, 차이 비교(diff),
버전 관리, 프로그램 처리에 그대로 쓸 수 있을 것.

\subsection{문법 정의}

\begin{table}[t]
\centering\small
\caption{선율 텍스트 문법 요약}
\label{tab:grammar}
\begin{tabularx}{\columnwidth}{@{}lX@{}}
\toprule
표기 & 의미 \\
\midrule
줄바꿈 & 각(세로열, 마디) 구분 \\
\code{|} & 정간(1박) 구분 \\
정간 안 공백 & 분박 --- 칸을 위$\to$아래로 등분 \\
공백 없는 연접 & 붙임 --- 한 분박 행 안의 가로 배치 \\
황·대·태$\cdots$응 & 12율 율명(한글 입력) \\
청·중청·배·하배 & 옥타브 접두어($+1,+2,-1,-2$) \\
\code{-} & 이음 --- 앞 음을 지속 \\
\code{쉼} & 쉼표(무음) \\
음 뒤 \code{<} & 숨표 --- 정간 오른쪽 아래 표시 \\
\code{\{이름\}} & 시김새 부가(예: \code{황\{미는표\}}) \\
\code{\{이름@$s$,$x$,$y$\}} & 시김새 크기 $s$\%·가로 $x$·세로 $y$ 조정 \\
\bottomrule
\end{tabularx}
\end{table}

표~\ref{tab:grammar}이 문법 전체다. 형식적으로 문서는 다음과 같이 정의된다.

\begin{quote}\small\ttfamily
문서\ \ := 각 (줄바꿈 각)* \\
각\ \ \ \ := 정간 ("|" 정간)* \\
정간\ \ := 분박 (" " 분박)* \\
분박\ \ := 음표+ \\
음표\ \ := (옥타브? 율명 | "-" | "쉼") 시김새* "<"?
\end{quote}

핵심 규칙은 두 가지다. 첫째, 정간 안의 \emph{공백}이 칸을 세로로 등분하고
(2분박·3분박$\cdots$), 공백 없이 이어 쓴 음은 그 분박 행 안에서 가로로
배치된다. 예컨대 \code{황 태협}은 첫 행에 황, 둘째 행에 태·협이 놓인
3음 구조로, 위치가 곧 시가라는 정간보 원리가 텍스트에도 그대로 성립한다.
둘째, 격자는 입력된 \code{|}의 개수가 아니라 문서 설정(한 각의 정간 수)을
기준으로 항상 온전히 그려지므로, 빈 정간(\code{ | })은 빈 채 유지되어
입력 도중의 불완전한 텍스트도 언제나 유효한 악보로 렌더링된다.

가사와 장단(장구 구음)도 동일한 문법을 공유한다. 가사는 정간마다 한 음절을
적으면 해당 정간 오른쪽에 병기되고, 장단은 덩·기덕·궁·덕·더러러러 등
구음 7종을 같은 구분자 규칙으로 적어 곡 첫머리에 한 각 폭으로 표시된다.
세 계층(선율·장단·가사)은 각·정간 단위로 1:1 정렬되며, 각을 삽입·삭제하면
가사 계층이 함께 밀리고 당겨져 정렬이 깨지지 않는다.

\subsection{시김새의 텍스트 직렬화}

시김새는 이름 기반으로 \code{\{미는표\}}처럼 부가하며, GUI에서 개별 기호의
크기·위치를 미세 조정하면 그 값이 \code{\{미는표@120,5,-5\}}(120\% 크기,
오른쪽 5, 위 5) 형태로 같은 텍스트에 직렬화된다. 기본값(100,0,0)이면 이름만
남는다. 이로써 ``문법으로 표현 못 하는 시각 조정''이 생기지 않으며,
텍스트 하나가 문서의 완전한 원천(single source of truth)이 된다.

\subsection{XML 기술 방식과의 간결성 비교}

기존 XML DTD 방식~\cite{lee2013xml}에서 음 하나는 최소한
여는 태그·속성·닫는 태그를 요구한다. 예를 들어 한 정간에서
황과 태를 2분박으로 나누는 경우를 개념적으로 비교하면 다음과 같다.

\begin{quote}\small
제안 문법: \code{황 태} (3자)\\[2pt]
XML 방식: \code{<jeonggan><bunbak><note pitch="황"/>}\allowbreak
\code{</bunbak><bunbak><note pitch="태"/></bunbak></jeonggan>} (수십 자)
\end{quote}

\TODO{lee2013xml 원문의 실제 태그 구조를 확인해 위 예시를 원문 표기로
교체하고, 취타 1곡 전체를 양 방식으로 인코딩한 문자 수·구문 요소 수를
표로 제시할 것}

물론 XML은 스키마 검증·범용 파서라는 장점이 있으므로, 제안 시스템은
문서 교환 시 텍스트 문법을 JSON 컨테이너에 담는 절충을 취했다(4.8절).
